% 木星（综述）
% license CCBYSA3
% type Wiki

本文根据 CC-BY-SA 协议转载翻译自维基百科\href{https://en.wikipedia.org/wiki/Jupiter}{相关文章}。


木星是距离太阳第五颗行星，也是太阳系中体积最大的行星。它是一颗气态巨行星，其质量几乎是太阳系中所有其他行星总和的 2.5 倍，但仍不到太阳质量的千分之一。它的直径是地球的 11 倍，是太阳的十分之一。木星以 5.20 天文单位（778.5 吉米）的平均距离绕太阳运行，公转周期为 11.86 年。在地球夜空中，木星是仅次于月球和金星的第三亮的天然天体，自史前时代以来便为人类所观察。它的名字源自古罗马宗教中至高神朱庇特（Jupiter）。

木星是太阳行星中最先形成的一颗，其在太阳系原始阶段向内迁移的过程深刻影响了其他行星的形成历史。木星的大气按质量计由 76\% 的氢和 24\% 的氦组成，内部密度更高。它还含有微量的碳、氧、硫、氖等元素，以及氨、水蒸气、磷化氢、硫化氢和烃类等化合物。木星的氦含量约为太阳的 80\%，与土星的成分相似。

木星外层大气被分为一系列纬向云带，这些云带在交界处产生湍流和风暴；其中最显著的结果就是“大红斑”——自 1831 年以来被记录的一场巨型风暴。由于木星自转极快（约 10 小时自转一周），使其呈现为一颗扁球体；与两极相比，赤道部位存在约 6.5\% 的明显隆起。科学家认为，木星内部结构包括一层由流体金属氢组成的外地幔，以及由更高密度物质构成的弥散内核。木星内部持续的收缩过程释放出比其从太阳接收的还要多的热量。木星的磁场是太阳系中最强且连续范围第二大的结构，由流体金属氢内部的涡电流产生。太阳风与其磁层相互作用，使磁层向外延伸并影响木星的运行环境。

至少有 97 颗卫星绕木星运行；其中四颗最大的卫星——木卫一（Io）、木卫二（Europa）、木卫三（Ganymede）和木卫四（Callisto）——位于其磁层内，用普通双筒望远镜即可观测。四者中最大的木卫三比行星水星还要大。木星周围还存在一个暗淡的行星环系统。木星环主要由尘埃组成，分为三部分：内侧呈光环状尘粒环的“晕环”（halo）、相对明亮的主环（main ring），以及外侧的薄雾状“薄环”（gossamer ring）。这些行星环在可见光和近红外光下呈现红色。木星环系统的年龄尚不清楚，可能追溯至木星形成时期。自 1973 年以来，已有九个机器人探测器访问过木星：包括七次飞掠任务和两艘专门的轨道探测器（另有两艘在途中）。在其他行星系统中也已发现类似木星的系外行星。
